\chapter{Análisis y diseño}

\section{Computabilidad}

El problema de crear un chat que utilice el protocolo indicado es computable, ya que consiste en la creación de dos programas (cliente y servidor) que utilizan algoritmos con un número finito de pasos, los cuales son las 12 funciones que describe el protocolo, que a su vez se descomponen en algoritmos con un número finito de pasos también.

\section{Comunicación}
La comunicación entre el cliente y el servidor se realiza mediante el uso de \textit{sockets} y el protocolo TCP/IP, donde el servidor conecta su socket a su dirección IP local en algún puerto y comienza a escuchar y, posteriormente, múltiples clientes conectan sus sockets a la IP y puerto del servidor, lo cual establece un canal de comunicación entre ambos. Una vez conectados, la comunicación comienza con el usuario, haciendo uso de la interfaz. para mandar un mensaje al servidor desde su programa, pasando por los sockets, llegando al servidor, decodificando e interpretando el mensaje, para posteriormente ejecutar acciones que cambien los registros de datos dentro del servidor y responder al mensaje enviado o notificar al resto de clientes conectados sobre la modificación.
\\

Por el lado del cliente, hay dos flujos de entrada y salida: uno de envío y otro de recepción de mensajes. A continuación, un ejemplo de ambos:
\begin{enumerate}
	\item Flujo de envío: Si el programa del cliente ya está en comunicación con el servidor, el usuario ingresa su nombre mediante una interfaz (entrada), esta valida que el formato sea correcto y luego se le notifica al cliente que envíe un mensaje que al servidor para identificarse. El flujo termina cuando cuando el mensaje es enviado (salida).
	\item Flujo de recepción: El cliente se encuentra esperando una respuesta del servidor y recibe una que indica que la identificación del usuario fue exitosa (entrada), entonces notifica a la interfaz para que le muestre el chat (salida).
\end{enumerate}

Por el lado del servidor hay un único flujo de entrada y múltiples salidas. A continuación un ejemplo:
\begin{enumerate}
	\item El servidor se encuentra esperando recibir un mensaje y en cuanto lo recibe valida que tenga el formato correcto y lo interpreta (entrada). Este mensaje es una solicitud de un cliente para identificar a un usuario, entonces valida que el formato del nombre sea correcto o si ya existe en los registros y si no, entonces lo registra y le envía un mensaje de respuesta diciendo que la operación fue exitosa (salida), además de notificar al resto de usuarios ya conectados que hay un nuevo usuario en el chat (múltiples salidas).
\end{enumerate}

La unión de estos flujos (los dos del cliente y el del servidor) da como resultado mínimo una entrada desde el usuario y una salida (posiblemente vacía) al mismo usuario, así que es suficiente para manejar al menos a un cliente; pero además, al esperar entradas y tener múltiples salidas (el servidor) se completa y resuelve el problema de mantener múltiples usuarios, dando como resultado la solución al problema de crear el chat, y por finalizado el análisis. Las posteriores decisiones se fundamentan en este análisis.

Más adelante se expone el motivo del paradigma escogido; sin embargo, a partir de ahora el diseño del cliente y servidor se inclinan a la orientación a objetos. Así mismo, se indica al lector que en el repositorio público de este proyecto se encuentran los diagramas necesarios para el diseño e implementación, así que aquí solamente se expondrá la parte teórica.

\section{Diseño del servidor}

El servidor funciona mediante las siguientes entidades y eventos:
\begin{itemize}
	\item \textbf{Entidades que contienen y validan datos:} Son las entidades que almacenan datos (como nombres de usuario, estados, nombres de sala, miembros, invitaciones, etc.) y que validan que su estado sea el adecuado. Entre estas entidades se encuentran un \textit{usuario}, una \textit{sala} y una entidad que valida que el formato de textos enviados a salas y usuarios sea el adecuado.
	\item \textbf{Encargados de la conexión:} Son entidades que se encargan de la comunicación directa con los sockets de los clientes para la escritura o lectura de mensajes.
	\item \textbf{Notificaciones de entrada/salida de mensajes:} para la entrada y la salida de mensajes del servidor.
	\item \textbf{Contenedor global de datos:} Contiene la información concreta de usuarios y salas existentes en el servidor y permite operar con ella.
	\item \textbf{Actores en el servidor:} Modifican el estado interno del servidor (sus datos, el estado de sus entidades, etc.) como respuesta a un mensaje recibido; además, tienen salida a otros clientes (manejan notificaciones).
	\item \textbf{Traductores de mensajes:} Traducen los mensajes recibidos y emitidos por el servidor a un formato apto para la transmisión de datos de sockets o a actores dentro del servidor.
\end{itemize}
Estas entidades se dividen, a su vez, en varios objetos implementados en el servidor y algunos mostrados en los diagramas de clases.

\section{Diseño del cliente}

El cliente funciona por medio de las siguientes entidades y eventos:
\begin{itemize}
	\item \textbf{Entidades observadas:} Son entidades que almacenan datos (nombres de usuario, estados, etc.) y los validan, pero cuyas modificaciones (cambio de estado) pueden ser notadas por las entidades que las utilicen.
	\item \textbf{Conexiones con el servidor:} Son entidades encargadas de realizar la escritura y lectura directa de mensajes con el servidor.
	\item \textbf{Emisores y receptores de mensajes:} Sirven como puente o interfaz entre las acciones que puede realizar el usuario (descritas en el protocolo) y la emisión de dicho mensaje o la recepción de alguno, mediante las conexiones con el servidor.
	\item \textbf{Traductores de mensajes:} Sirven para traducir mensajes recibidos y emitidos por el cliente al servidor en un formato compatible con los sockets.
	\item \textbf{Proveedores de contexto:} Son todas las entidades que proveen datos y comandos como puente entre el usuario y la emisión/recepción de mensajes.
	\item \textbf{Vistas:} Son las entidades encargadas de comunicarse con los proveedores de contexto para comunicarse a su vez con el servidor. Dependen del contexto para saber qué mostrarle y permitirle hacer al usuario.
\end{itemize}
